\subsection{QG9-M1: windowed comparability and an auditable error budget (protocol vs.\ finite-order EFT)}
\label{app:qg9_m1_windowed_comparability}

\paragraph{Scope.}
\AuditTag This appendix instantiates the generic theoremized error budget of Appendix~\ref{app:eft_error_bounds} in the specific
``QG interface'' context used by this manuscript:
protocol-side observables produced by the deterministic pipelines in Section~\ref{sec:qg_interface_full_fusion}
are compared to the declared finite-order EFT representative under explicitly scoped matching conventions.
This is a \emph{windowed} statement: it does not claim a theorem-level UV completion.

\subsubsection{Declared UV window and protocol state}
\label{subsec:qg9_m1_window_state}

\begin{definition}[UV window in the resolution coordinate]\label{def:qg9_uv_window}
\InterfaceTag Fix a resolution coordinate $r$ and a reference map $\mu(r)=\mu_0\varphi^{\,r}$ (Appendix~\ref{app:running_couplings_resolution_flow}).
A \emph{UV window} is a bounded interval
$
W=[r_{\min},r_{\max}]
$
equivalently a bounded scale interval
$
[\mu_{\min},\mu_{\max}]=[\mu(r_{\min}),\mu(r_{\max})].
$
All comparability statements in this appendix are conditional on working within this declared window.
\end{definition}

\begin{definition}[Declared protocol state]\label{def:qg9_protocol_state}
\InterfaceTag Fix a protocol state $(m,n,K)$ as in Appendix~\ref{subsec:tick_cap_joint_key_contract}.
All protocol-side estimators in the registry below are evaluated at this fixed $(m,n,K)$, and any counterfactual sweeps
(kernel-family sweeps, smoothing sweeps, scheme sweeps) are restricted to explicit finite families declared by the manuscript.
\end{definition}

\subsubsection{A concrete finite observable registry for the QG-interface suite}
\label{subsec:qg9_m1_registry}

\paragraph{Default contract instance (generated; reproducibility-first).}
\AuditTag To avoid underspecification, we record a compact reference instance of the QG9-M1 contract (a UV window expression and an anchor protocol state)
as a reproducibility aid. It is generated deterministically from the protocol-state selection artifacts by
\texttt{scripts/exp\_qg9\_windowed\_comparability\_pack.py}:
\noindent\AuditTag \textbf{Reference QG9-M1 instance (reproducibility).} The following is a non-normative instance used to pin down a concrete window/state when reproducing the QG-interface comparability tables; it does not add assumptions beyond the declared staircase dictionary.
\begin{itemize}
\item \textbf{Protocol state anchor:} \texttt{gamma_direct} at $\mu=$ SPARC-direct, $(m,n)=(6,3)$, $K=$ \texttt{tempered_deg}:0 (from \texttt{sections/generated/protocol\_state\_selected.json}).
\item \textbf{Reference UV window:} use the staircase threshold coordinate $r_{\mathrm{th}}(m)=(m-6)\,r_{\mathrm{step}}$ with $r_{\mathrm{step}}=2\pi$ (Section~\ref{subsec:p3_resolution_jumps}) and take $W_{\mathrm{ref}}=[r_{\mathrm{th}}(6),\,r_{\mathrm{th}}(10)]=[0,\,8\pi]$ as a minimal even-step window aligned with the QG interface suite family $(m,n)\in\{(6,3),(8,4),(10,5)\}$.
\item \textbf{EFT order and confidence:} take $N=1$ and $\delta=10^{-2}$ as a compact default for the registry-wide bound in Appendix~\ref{app:eft_error_bounds} (union-bound over a finite registry).
\end{itemize}

\paragraph{Registry instance (generated; audit index).}
\AuditTag We record a compact concrete choice of $\mathfrak{Obs}_{\le N}$ (Definition~\ref{def:obs_registry_finite}) tailored to the QG-interface suite
and to the weak-field/time-dictionary carriers emphasized in Section~\ref{sec:falsifiability}.
The registry is generated deterministically by \texttt{scripts/exp\_qg9\_windowed\_comparability\_registry.py}:
\texttt{scripts/exp\_qg9\_windowed\_comparability\_pack.py}:
\begin{center}
\scriptsize
\setlength{\tabcolsep}{6pt}
\renewcommand{\arraystretch}{1.15}
\begin{tabular}{p{0.22\textwidth} p{0.25\textwidth} p{0.23\textwidth} p{0.24\textwidth}}
\toprule
observable $\mathcal{O}$ & protocol-side rule & EFT-side rule/scope & error-budget terms \\
\midrule
$\widehat\chi$ (reconstructed overhead proxy field) & Deterministic reconstruction pipeline (Appendix~\ref{app:overhead_to_gravity_closure}; Appendix~\ref{app:protocol_to_continuum_error_control}). & EFT-side target is the corresponding scalar proxy under the declared representative action and matching conventions (Appendix~\ref{app:cap_continuum_action_closure}). & $E_{\mathrm{stat}}$+$E_{\mathrm{disc}}$ from Appendix~\ref{app:protocol_to_continuum_error_control}; $E_{\mathrm{bias}}$ from declared kernel/smoothing sweeps. \\
$\widehat G_{00,h}$ (weak-field curvature proxy from $\widehat\chi$) & $\widehat G_{00,h}=-2\widehat\gamma\,\Delta_h\widehat\chi_h$ (Appendix~\ref{app:weak_field_curvature_from_chi}). & Weak-field EFT prediction in the declared truncation scope (Appendix~\ref{app:variational_field_equations}; Appendix~\ref{app:overhead_to_gravity_closure}). & $E_{\mathrm{disc}}$ via Laplacian truncation/noise amplification; $E_{\mathrm{match}}$ includes $\gamma$-dictionary matching envelope when applicable. \\
$\tau_{\mathrm{WS}}(\omega)$ / $N_{\mathrm{WS}}$ (delay-derived lapse proxy) & Wigner--Smith delay dictionary and audit gates (Appendix~\ref{app:time_mass_delay}; Section~\ref{subsec:p6_wigner_smith_delay}). & EFT-side target specified by the declared time dictionary and carrier conventions (Appendix~\ref{app:renormalization_dictionary_and_boundaries}). & $E_{\mathrm{match}}$ over a finite carrier/scheme family; endpoint/unwrap failures trigger explicit failure points (S1). \\
$|\mathcal{R}_\star|/4^n$ (budget-triggered horizon occupancy fraction) & Capacity-only budget trigger applied to reconstructed $\widehat\chi$ (Section~\ref{subsec:p8_chi_horizon_budget}; Appendix~\ref{app:protocol_horizon_tick_trap}). & EFT-side interpretation is interface-level only; comparability uses the declared capacity/entropy dictionary (Appendix~\ref{app:thermodynamics_from_equivalence}). & $E_{\mathrm{bias}}$ includes the declared margin constant and smoothing; $E_{\mathrm{stat}}$ inherits from $\widehat\chi$ uncertainty. \\
Full-fusion interface gates (ledger closure; $V^2+D^2\le 1$; counterfactual deltas) & Deterministic full-fusion generator and artifacts (Section~\ref{sec:qg_interface_full_fusion}). & No EFT-side target: these are protocol-interface consistency constraints (Prediction~P9; Section~\ref{subsec:p9_full_fusion_interface_gates}). & Must hold exactly up to numerical tolerance; violations falsify the interface contract (not the theorem-level core). \\
\bottomrule
\end{tabular}
\end{center}

\subsubsection{Windowed comparability contract and the error budget \texorpdfstring{$\epsilon_N$}{epsilonN}}
\label{subsec:qg9_m1_contract}

\paragraph{Contract form.}
\MathTag Fix $N\ge 1$ and confidence level $\delta\in(0,1)$.
Within a declared UV window $W$ (Definition~\ref{def:qg9_uv_window}) and at a fixed protocol state $(m,n,K)$
(Definition~\ref{def:qg9_protocol_state}), the QG9-M1 target is a registry-wide bound of the form
$$
\forall\,\mathcal{O}\in\mathfrak{Obs}_{\le N},\qquad
\left|\langle\mathcal{O}\rangle_{\mathrm{prot}}-\langle\mathcal{O}\rangle_{\mathrm{EFT}\le N}\right|
\le \epsilon_N(m,n;\delta),
$$
where the right-hand side is an \emph{auditable} quantity assembled from explicit components (Appendix~\ref{app:eft_error_bounds}).

\begin{remark}[Where $K$ enters]\label{rem:qg9_kernel_in_epsilon}
\AuditTag The generic bound in Appendix~\ref{app:eft_error_bounds} is written as $\epsilon_N(m,n;\delta)$ for brevity.
In QG-interface applications, $K$ enters through the registry specification (Definition~\ref{def:obs_registry_finite}):
it fixes the protocol-side estimator and can be included in bias/sensitivity terms.
When emphasizing this conditioning, one may write $\epsilon_N(m,n,K;\delta)$ without changing the logic of the decomposition.
\end{remark}

\paragraph{Explicit components and provenance.}
\AuditTag For the QG-interface observables listed in the registry, the five-term decomposition of Appendix~\ref{app:eft_error_bounds}
is intended to be instantiated as follows:
\begin{itemize}
\item \textbf{Protocol statistical and discretization components ($E_{\mathrm{stat}}$, $E_{\mathrm{disc}}$).}
  Use the protocol-to-continuum error control module (Appendix~\ref{app:protocol_to_continuum_error_control})
  for $\widehat\chi$-based channels and their derivative descendants, and the weak-field curvature-from-$\chi$ module
  (Appendix~\ref{app:weak_field_curvature_from_chi}) for the $\widehat G_{00}$ proxy chain.
\item \textbf{Matching/scheme envelope ($E_{\mathrm{match}}$).}
  Restrict all scheme/threshold conventions to an explicit finite family and budget its width as an envelope
  (Appendix~\ref{app:matching_envelope_theoremization}) under the scheme-invariance contract (Appendix~\ref{app:scheme_invariance_audit_contract}).
\item \textbf{Truncation remainder ($E_{\mathrm{trunc}}$).}
  When an EFT-side prediction is invoked, the truncation scope and remainder model must be declared explicitly;
  if an EG-style remainder envelope is assumed, it can be integrated via Appendix~\ref{app:strong_eft_remainder_bounds}.
\item \textbf{Protocol bias and counterfactual sensitivity ($E_{\mathrm{bias}}$).}
  Bias terms must be tracked as explicit finite-family counterfactuals (kernel-family, smoothing, estimator choices),
  and reported as envelopes rather than treated as implicit knobs (cf.\ Appendix~\ref{app:global_model_selection_mdl} for
  look-elsewhere accounting across families).
\end{itemize}

\paragraph{Budget mapping table (generated; audit index).}
\AuditTag The following table records, for each registry item, where each budget component is sourced from
(hard bounds, envelopes, or explicit counterfactual sweeps). It is generated deterministically by
\texttt{scripts/exp\_qg9\_windowed\_comparability\_pack.py}:
\begin{center}
\scriptsize
\setlength{\tabcolsep}{6pt}
\renewcommand{\arraystretch}{1.15}
\begin{tabular}{p{0.12\textwidth} p{0.26\textwidth} p{0.20\textwidth} p{0.18\textwidth} p{0.18\textwidth}}
\toprule
$\mathcal{O}$ & $E_{\mathrm{stat}}/E_{\mathrm{disc}}$ provenance & $E_{\mathrm{match}}$ provenance & $E_{\mathrm{trunc}}$ scope & $E_{\mathrm{bias}}$ provenance \\
\midrule
$\widehat\chi$ & Appendix~\ref{app:protocol_to_continuum_error_control} (concentration + log propagation + derivative tradeoffs). & None unless a matching dictionary is invoked for units/anchors (Match scope). & Not applicable unless an EFT-side target is specified at finite $D_{\max}$. & Kernel-family / smoothing / estimator sweeps (explicit finite families; audit tables). \\
$\widehat G_{00,h}$ & Appendix~\ref{app:weak_field_curvature_from_chi} (discrete proxy + error budget; Laplacian truncation/noise amp). & $\gamma$ dictionary matching and scheme envelope when applicable (Appendix~\ref{app:matching_envelope_theoremization}; Appendix~\ref{app:scheme_invariance_audit_contract}). & Only if compared to an EFT-side weak-field prediction; then use Appendix~\ref{app:eft_error_bounds}. & Reconstruction/regularization sweeps (bounded counterfactuals). \\
$\tau_{\mathrm{WS}}(\omega)$ / $N_{\mathrm{WS}}$ & Numerical differentiation/unwrap stability gates (Appendix~\ref{app:time_mass_delay}; Appendix~\ref{app:force_phase_delay_audit}). & Finite carrier/scheme family envelope (Appendix~\ref{app:matching_envelope_theoremization}). & EFT-side truncation only if an EFT prediction is used; otherwise this remains an interface dictionary. & Observation-class boundary: endpoint-gated/intermittent channels (Section~\ref{subsec:p9_full_fusion_interface_gates}). \\
$|\mathcal{R}_\star|/4^n$ & Inherits from $\widehat\chi$ uncertainty; capacity-only computation is exact given $(m,n)$. & None (capacity-only); any physical reinterpretation is Match/Iface only. & Not applicable (no EFT-side target claimed). & Margin constant and smoothing rule must be fixed or swept as an explicit finite family. \\
Full-fusion gates & Exact interface checks (ledger closure; complementarity); numerical tolerance only. & None (protocol-interface consistency only). & Not applicable. & Counterfactual baseline required (wormhole off) and observation-class limitations must be stated. \\
\bottomrule
\end{tabular}
\end{center}

\paragraph{Evidence pointer table (generated; audit index).}
\AuditTag The following table binds each registry item to concrete reproducible artifacts in \texttt{sections/generated/}
and the generator scripts that produce them. It is generated deterministically by
\texttt{scripts/exp\_qg9\_windowed\_comparability\_pack.py}:
\begin{center}
\scriptsize
\setlength{\tabcolsep}{6pt}
\renewcommand{\arraystretch}{1.15}
\begin{tabular}{p{0.12\textwidth} p{0.40\textwidth} p{0.22\textwidth} p{0.20\textwidth}}
\toprule
$\mathcal{O}$ & reproducible artifacts & generator scripts & audit note \\
\midrule
$\widehat\chi$ & \texttt{sections/generated/qg\_interface\_suite\_rows.tex}; \texttt{sections/generated/qg\_interface\_suite\_summary.tex}; \texttt{sections/generated/chi\_kernel\_family\_sweep\_rows.tex}. & \texttt{scripts/exp\_qg\_interface\_suite.py}; \texttt{scripts/exp\_chi\_kernel\_family\_sweep.py}. & Counterfactual sensitivity of reconstructed scalars is recorded via kernel-family sweeps. \\
$\widehat G_{00,h}$ & \texttt{sections/generated/curvature\_e2e\_rows.tex}; \texttt{sections/generated/curvature\_e2e\_summary.tex}; \texttt{sections/generated/curvature\_e2e\_gamma\_rows.tex}; \texttt{sections/generated/curvature\_e2e\_gamma\_summary.tex}; \texttt{sections/generated/curvature\_e2e\_gamma\_stability\_rows.tex}. & \texttt{scripts/exp\_curvature\_bridge\_end\_to\_end.py}. & End-to-end provenance: protocolized input $\rightarrow \widehat\chi \rightarrow$ curvature proxy, including $\gamma$ diagnostics and bounded stability rows. \\
$\tau_{\mathrm{WS}}(\omega)$ / $N_{\mathrm{WS}}$ & \texttt{sections/generated/qg\_interface\_suite\_rows.tex}; \texttt{sections/generated/qg\_interface\_suite\_summary.tex}; \texttt{sections/generated/force\_phase\_delay\_audit\_toy\_rows.tex}. & \texttt{scripts/exp\_qg\_interface\_suite.py}; \texttt{scripts/exp\_force\_phase\_delay\_audit\_toy.py}. & Numerical stability of phase$\rightarrow$delay differentiation is recorded as an auditable toy table; carriers must satisfy S1-style gates. \\
$|\mathcal{R}_\star|/4^n$ & \texttt{sections/generated/chi\_horizon\_budget\_occupancy\_rows.tex}; \texttt{sections/generated/chi\_horizon\_budget\_occupancy\_summary.tex}. & \texttt{scripts/exp\_chi\_horizon\_budget\_occupancy.py}. & Capacity-only occupancy is deterministic given $(m,n)$ and the declared budget/margin. \\
Full-fusion gates & \texttt{sections/generated/full\_fusion\_rows.tex}; \texttt{sections/generated/full\_fusion\_nowh\_rows.tex}; \texttt{sections/generated/full\_fusion\_compare\_rows.tex}; \texttt{sections/generated/full\_fusion\_summary.tex}; \texttt{sections/generated/full\_fusion\_wormhole\_sweep\_rows.tex}; \texttt{sections/generated/full\_fusion\_wormhole\_pareto\_rows.tex}. & \texttt{scripts/exp\_full\_fusion\_bh\_wormhole\_measurement.py}; \texttt{scripts/exp\_full\_fusion\_wormhole\_sweep.py}; \texttt{scripts/exp\_full\_fusion\_wormhole\_adaptive\_search.py}. & Gate checks are ledger closure + complementarity + counterfactual isolation; sweep outputs provide bounded trade-off records. \\
\bottomrule
\end{tabular}
\end{center}

\paragraph{Minimal acceptance checklist (generated; audit gate).}
\AuditTag The following checklist summarizes the minimal gates required to treat QG9-M1 as an auditable deliverable:
it makes underspecification and hidden family enlargements explicit failure modes.
It is generated deterministically by \texttt{scripts/exp\_qg9\_windowed\_comparability\_pack.py}:
\begin{center}
\scriptsize
\setlength{\tabcolsep}{6pt}
\renewcommand{\arraystretch}{1.15}
\begin{tabular}{p{0.16\textwidth} p{0.40\textwidth} p{0.38\textwidth}}
\toprule
gate & acceptance requirement & evidence pointers \\
\midrule
E0 (default instance pinned) & A concrete reference instance (window + state + $N,\delta$) is recorded to avoid underspecification. & \texttt{sections/generated/qg9\_windowed\_comparability\_default\_instance.tex}. \\
E1 (finite registry) & $\mathfrak{Obs}_{\le N}$ is finite and explicitly listed; each $\mathcal{O}$ has protocol/EFT rules and a metric. & Appendix~\ref{app:eft_error_bounds} (Def.~\ref{def:obs_registry_finite}); \texttt{sections/generated/qg9\_windowed\_comparability\_registry.tex}. \\
E2 (estimator semantics) & Protocol-side estimators (smoothing, differencing, kernels) are declared or swept in explicit finite families. & \texttt{sections/generated/qg9\_windowed\_comparability\_budget.tex}; \texttt{sections/generated/qg9\_windowed\_comparability\_evidence.tex}. \\
E3 (truncation scope) & If an EFT-side target is invoked: $N$ and $D_{\max}$ (and any remainder model/envelope) are declared; otherwise the item is interface-only. & Appendix~\ref{app:eft_error_bounds} (E3); Appendix~\ref{app:strong_eft_remainder_bounds} when a remainder envelope is used. \\
E4 (scheme/threshold scope) & Matching conventions are fixed or bounded into an explicit finite family; envelope width is budgeted (no hidden scheme knobs). & Appendix~\ref{app:matching_envelope_theoremization}; Appendix~\ref{app:scheme_invariance_audit_contract}. \\
R1/R2 (finite-family envelopes) & All counterfactual families (kernels, schemes, carriers) are explicit and treated as envelopes, not free parameters. & \texttt{sections/generated/qg9\_windowed\_comparability\_budget.tex}; MDL registry discipline: Appendix~\ref{app:global_model_selection_mdl}. \\
Obs-class boundary & Observation-class identifiability limits are stated (e.g.\ endpoint-gated channels; final-snapshot deltas may miss intermittent effects). & Section~\ref{subsec:p9_full_fusion_interface_gates}; \texttt{sections/generated/full\_fusion\_compare\_rows.tex}. \\
Provenance gate & Artifacts are content-addressed for auditable provenance (script/deps fingerprint). & \texttt{sections/generated/artifact\_hash\_registry\_summary.tex}. \\
\bottomrule
\end{tabular}
\end{center}

\paragraph{Numerical instantiation (generated; actual computation).}
\AuditTag Finally, we record a concrete numerical instantiation of the QG9-M1 budget for the default reference instance:
the values and error envelopes are computed deterministically from the paper's carriers and generated artifacts (not hand-entered).
This fragment is generated by \texttt{scripts/exp\_qg9\_windowed\_comparability\_pack.py}:
\noindent\AuditTag QG9-M1 numerical instantiation (default instance): anchor $(m,n)=(6,3)$, $K=$ \texttt{tempered_deg}:0, $N=1$, $\delta=0.0100$. All quantities below are computed deterministically from the paper's generated artifacts/scripts.
\begin{center}
\scriptsize
\setlength{\tabcolsep}{6pt}
\renewcommand{\arraystretch}{1.15}
\begin{tabular}{p{0.22\textwidth} p{0.20\textwidth} p{0.20\textwidth} p{0.32\textwidth}}
\toprule
$\mathcal{O}$ & numeric proxy/value & $\epsilon$ (numeric) & audit note \\
\midrule
$\widehat\chi$ (sup error) & 0.356337 & 0.356337 & Computed as $\|\widehat\chi-\chi_{\rm true}\|_\infty$ in the end-to-end reconstruction carrier at $n=3,m=6$. \\
$\Delta_h\widehat\chi$ (Laplacian-stage bound) & 54.759153 & 183.459268 & Computed bound $C h^2+(4d/h^2)\epsilon_\chi$ and realized error $\|\Delta_h\widehat\chi-\Delta\chi_{\rm true}\|_\infty$. \\
$\widehat G_{00,h}$ (proxy) & 0.000068 & 0.000068 & Bound uses $\gamma$ family envelope (min/max/mean) and $\max|\Delta\chi_{\rm true}|$ under the same carrier. \\
$\tau_{\mathrm{WS}}$ (peak) & 4.000000 & 0.0 & Breit--Wigner one-channel benchmark in the QG interface suite (analytic carrier). \\
$N_{\mathrm{WS}}$ (peak/far) & 0.006897/1.000000 & 0.0 & Delay-to-lapse mapping in the same analytic carrier; numerical error is zero in this benchmark. \\
$|\mathcal{R}_\star|/4^n$ (occupancy) & 0.171875 & 0.0 & Capacity-only occupancy fraction for the first feasible row at the anchor $(m,n)$ in \texttt{qg\_interface\_suite\_rows.tex}. \\
Full-fusion ledger gate (max resid; on/off) & 0.000001/0.000001 & 0.000001 & Computed as $\max_t |E_{\rm tot}-(E_{\rm part}+E_{\rm field}+E_{\rm emit}+E_{\rm wh})|$ from generated time-series rows. \\
Full-fusion complementarity (max violation; free/trap) & 0.000000/0.000000 & 0.000000 & Computed as $\max_t \max(0,V^2+D^2-1)$ (wormhole on); off-baseline computed analogously. \\
\bottomrule
\end{tabular}
\end{center}

\paragraph{Minimal failure points (window discipline).}
\AuditTag This windowed comparability contract fails only if the registry/window is underspecified or if uncontrolled expansions of families are introduced.
The corresponding normalized failure-point templates are:
missing finite registry (E1), missing estimator semantics (E2), missing truncation scope (E3), and missing scheme scope (E4)
as recorded in Appendix~\ref{app:eft_error_bounds}, together with the cross-cutting failure-point registry in Appendix~\ref{app:minimal_failure_point_templates}.

